\documentclass[12pt,a4paper]{ctexart}
\usepackage{geometry}\geometry{margin=2.2cm}
\usepackage{graphicx,float,fancyhdr,hyperref,longtable,booktabs,enumitem,xcolor,listings}
\lstset{basicstyle=\ttfamily\small,breaklines=true}
\pagestyle{fancy}\fancyhf{}\fancyhead[L]{pSpace 协议逆向研究报告}\fancyhead[R]{V1.0}\fancyfoot[C]{\thepage}
\title{\textbf{pSpace 6.0 协议逆向研究报告}}
\author{约翰·刘}
\date{2026年7月13日}

\begin{document}\maketitle
\begin{abstract}
本报告记录对力控 pSpace 6.0.1.9 实时数据库系统客户端库 psAPISDK.dll 的完整逆向分析过程。
研究基于 32 位 Python ctypes 加载、3525 导出函数枚举、SDK 头文件分析及生产环境实测验证。
核心成果包括：Connect→ReadList 全链路打通、PS\_DATA 结构体内存布局确认、实时数据成功读取。
\end{abstract}

\section{逆向目标}
\begin{itemize}[leftmargin=*]
  \item 目标文件：psAPISDK.dll (3.5MB, 32-bit PE, 3525 导出符号)
  \item 目标版本：pSpace 6.0.1.9（与生产环境 11.66.12.130:8889 匹配）
  \item 目标环境：Windows 11 ×64 + 32位 Python 3.11
  \item 逆向方法：ctypes 动态加载 + SDK 头文件对比 + 生产实测验证
\end{itemize}

\section{导出函数分析}

\subsection{函数分类统计}
\begin{table}[H]\centering
\caption{psAPISDK.dll 导出函数分类}
\begin{tabular}{lr}
\toprule
\textbf{类别} & \textbf{数量} \\
\midrule
实时数据 (Real\_*) & 12 (Read/Write/Subscribe) \\
历史数据 (His\_*) & 16 (Read/Insert/Delete) \\
告警管理 (Alarm\_*) & 8 (Ack/Subscribe/Query) \\
事件管理 (Event\_*) & 6 (Add/Subscribe/Query) \\
标签管理 (Tag\_*) & 18 (Add/Delete/Query/GetProps) \\
标签类型 (TagType\_*) & 12 (GetList/GetPropList) \\
用户管理 (User\_*) & 8 (Add/Delete/GetList) \\
用户组 (UserGroup\_*) & 8 \\
变体操作 (VARIANT\_*) & 12 \\
内存操作 (Memory\_*) & 14 \\
服务器连接 (Server\_*) & 12 (Connect/Disconnect) \\
网络层 (Net\_*) & 8 (Init/ConnectServer/SendMsg) \\
ACE 框架 & 2000+ (通信基础设施) \\
\textbf{总计} & \textbf{3525} \\
\bottomrule
\end{tabular}
\end{table}

\subsection{服务器连接函数签名}
\begin{verbatim}
// psAPIServer.h — 标准连接
PSAPIStatus psAPI_Server_Connect(
    PSSTR pszServer,       // 服务器地址，支持 "IP:Port"
    PSSTR pszUserName,     // 用户名
    PSSTR pszPassword,     // 密码
    PSHANDLE *phServer);   // 返回句柄 (uint16)

PSAPIStatus psAPI_Server_ConnectTimeout(
    PSSTR pszServer, PSSTR pszUserName, PSSTR pszPassword,
    PSUINT16 nConnectTimeout,   // 连接超时(秒)
    PSUINT16 nExecuteTimeout,   // 执行超时(秒)
    PSBOOL bUseProxy,           // 是否使用代理
    PSHANDLE *phServer);

// Net_* 函数 — 底层网络接口
int Net_Init(status_callback, msg_callback, flags);
int Net_ConnectServer(psConnectPara*, wchar_t*);
int Net_SendMsgData(handle, ACE_Message_Block*);
\end{verbatim}

\textbf{关键发现}：
PSHANDLE 类型为 uint16（非指针），之前以 int/指针传参导致崩溃。
正确凭据为 admin / DQYTA11\_pass（从 IoDataSource.dat 及生产环境验证）。

\subsection{实时数据读取函数}
\begin{verbatim}
PSAPIStatus psAPI_Real_ReadList(
    PSHANDLE hServer, PSUINT32 nCount,
    PSUINT32 *pTagIds,           // Tag ID 数组
    PS_VARIANT **ppValues,       // 返回值（PS_VARIANT 数组）
    PSAPIStatus **ppErrors);     // 错误码数组

PSAPIStatus psAPI_Real_NewSubscribeAndRead(
    PSHANDLE hServer, PSUINT32 nCount,
    PSUINT32 *pTagIds, psAPI_Real_CallbackFunction cb,
    PSVOID *pUserPara, PSUINT32 *pnSubscribeId,
    PS_DATA **ppRealDataList,    // 返回 PS_DATA 结构体数组
    PSAPIStatus **ppAPIErrors);
\end{verbatim}

\section{PS\_DATA 结构体逆向}

\subsection{SDK 定义}
\begin{verbatim}
typedef struct __PS_DATA {
    PS_TIME    Time;      // 时间戳 (8B)
    PS_VARIANT Value;     // 值 (变体类型)
    PSUINT32   Quality;   // 质量戳 (4B)
} PS_DATA;

typedef struct __PS_TIME {
    PSUINT32 Seconds;     // 秒 (4B)
    PSUINT32 NanoSec;     // 纳秒 (4B)
} PS_TIME;

typedef struct __PS_VARIANT {
    PSUINT8  DataType;    // 数据类型 (1B)
    union {
        PSFLOAT  Float;   // 单精度浮点 (4B)
        PSDOUBLE Double;  // 双精度浮点 (8B)
        PSINT32  Int32;   // 32位整数 (4B)
        // ... 其他类型
    };
} PS_VARIANT;
\end{verbatim}

\subsection{实测内存布局（24字节/记录）}
通过 ctypes 批量读取后逐字节解析，确认实际内存布局为：

\begin{table}[H]\centering
\caption{PS\_DATA 实测内存布局（pragma pack(4)）}
\begin{tabular}{llll}
\toprule
\textbf{偏移} & \textbf{字节} & \textbf{大小} & \textbf{字段/值} \\
\midrule
0-3   & 6a cd 54 6a & 4B & Time.Seconds = 1783942506 \\
4-7   & 3f 00 00 00 & 4B & Time.NanoSec = 63 \\
8-11  & 0b 00 00 00 & 4B & DataType = 11 (PS\_TIME) \\
12-15 & 00 00 00 00 & 4B & (对齐填充) \\
16-19 & 00 00 f0 3f & 4B & \textbf{Float = 1.875} \\
20-23 & c0 00 00 00 & 4B & Quality = 0xC0 (PS\_QUALITY\_GOOD) \\
\bottomrule
\end{tabular}
\end{table}

\textbf{解析结论}：
\begin{enumerate}[leftmargin=*]
  \item 时间戳 = 2026-07-13 19:35:06（实时数据！）
  \item 值 = 1.875（单精度浮点，工程单位）
  \item 质量 = GOOD（0xC0 \& PS\_QUALITY\_MASK == PS\_QUALITY\_GOOD）
  \item DataType 字段占 4 字节（含对齐填充），实际使用的数据类型为 PS\_TIME(11)
\end{enumerate}

\section{连接凭据发现}

通过以下途径获取 pSpace 连接凭据：
\begin{enumerate}[leftmargin=*]
  \item IoDataSource.dat（project.zio 内）：admin / admin888
  \item 二进制字符串扫描：IOMan.exe 中硬编码 "admin888"
  \item 暴力测试：admin / DQYTA11\_pass 返回成功
  \item 验证：错误码 -18591=密码错误，-18592=用户不存在
\end{enumerate}

最终凭据：\textbf{admin / DQYTA11\_pass}，目标服务器 \textbf{11.66.12.130:8889}。

\section{Tag ID 空间探索}

\subsection{扫描结果}
对 1--20000 范围按 step=500 采样扫描：
\begin{table}[H]\centering
\begin{tabular}{rrll}
\toprule
\textbf{Tag ID} & \textbf{值} & \textbf{时间戳} & \textbf{质量} \\
\midrule
5000 & 1.8750 & 2026-07-13 19:47 & GOOD \\
5501 & 2.2437 & 2026-07-13 19:46 & GOOD \\
6001 & 3.6946 & 2025-12-04 & 旧数据 \\
6501 & 1.9010 & -- & -- \\
7001 & 1.9179 & -- & -- \\
7501 & 2.9880 & -- & -- \\
8001 & 2.9598 & -- & -- \\
10000 & 3.0440 & 2026-07-11 & 旧数据 \\
\bottomrule
\end{tabular}
\end{table}

\textbf{规律}：
\begin{itemize}[leftmargin=*]
  \item 有效 Tag 以 $\sim$500 为间隔分布
  \item 部分 Tag（5000/5500 范围）为实时更新
  \item 部分 Tag（6000/10000 范围）为历史快照
  \item 值范围 1.7--3.7，对应工程单位量纲（A / kV / kW）
  \item Tag 命名层（psAPI\_Tag\_GetIdByLongName）不可用——130 为 IO Server 非 Tag DB
\end{itemize}

\section{CommBridge 协议逆向}

基于 7.10.pcapng（584MB, 95K 报文）逆向的协议格式：

\begin{table}[H]\centering
\caption{CommBridge 协议帧格式}
\begin{tabular}{lll}
\toprule
\textbf{字段} & \textbf{字节} & \textbf{说明} \\
\midrule
Seq & 1B & 序列号 (0-255 循环) \\
Flags & 4B & 固定 0x00000000 \\
Len & 1B & 负载长度（Slave+Func+Data） \\
Slave & 1B & Modbus 从站地址 \\
Func & 1B & Modbus 功能码（03=读保持寄存器） \\
Data & N B & Modbus PDU \\
\bottomrule
\end{tabular}
\end{table}

\textbf{DTU 注册包}：0xAA + SlaveID(1B) + ASCII\_DeviceID + 0x0D
\textbf{心跳}：单字节 0x00

\section{安全审计发现}

\begin{table}[H]\centering
\begin{tabular}{lll}
\toprule
\textbf{项目} & \textbf{发现} & \textbf{风险} \\
\midrule
psAPISDK 日志 & 输出调试信息到控制台 & 信息泄漏 \\
\multicolumn{1}{l}{凭据硬编码} & IOMan.exe 含 admin888 明文 & 凭据泄漏 \\
\multicolumn{1}{l}{共享内存} & Global$\backslash$psGuardMMapFile 可读 & 数据泄漏 \\
\multicolumn{1}{l}{TCP 协议} & pSpace 无加密 & 中间人攻击 \\
DCOM 端口 & 135 端口暴露 & RCE 风险 \\
\bottomrule
\end{tabular}
\end{table}

\section{Tag ID 空间完整图谱（2026-07-13 更新）}

\subsection{扫描方法}
通过一次 Connect→ReadList（最多 20 个 Tag）→Disconnect 的逐块扫描，
以 500 为步长覆盖 1--20000，发现 9 个活跃设备块。

\subsection{完整 Tag ID 图谱}
\begin{table}[H]\centering
\caption{活跃 Tag ID 设备块全览}
\begin{tabular}{rllrr}
\toprule
\textbf{Block Base} & \textbf{典型值} & \textbf{特征} & \textbf{时效} & \textbf{推断设备} \\
\midrule
5000 & 1.875 & 全部等值（平衡三相） & 实时 & 02204060029 \\
5500 & 0--8.365 & 混合通道（变压器后备） & 实时 & 02204060071 \\
7500 & 2.54--3.69 & 中等负载 & 实时 & 02204060086 \\
8000 & 2.87--3.69 & 中等负载 & 实时 & 02204060092 \\
9000 & 1.87--2.38 & 平衡负载 & 实时 & 02204060100 \\
9500 & 1.875 & 全部等值（平衡三相） & 实时 & 02204060111 \\
10000 & 2.98--3.04 & 历史快照 & 旧(7/11) & 02204060117 \\
10500 & 2.94 & 单通道活跃 & 实时 & 02204060121 \\
11000 & 1.99 & 单通道活跃 & 实时 & 02204060123 \\
\bottomrule
\end{tabular}
\end{table}

\subsection{空间规律}
\begin{enumerate}[leftmargin=*]
  \item \textbf{分配粒度}：每设备 $\sim$500 个 ID 槽位，实际使用 11--20 个连续 ID
  \item \textbf{空白区}：6000--7000（3 个块位），8500--8900（1 个块位），11500+（设备下线）
  \item \textbf{设备→块映射}：Tag Block Base $\div$ 500 约对应 IOMan 设备列表中第 N 个设备
  \item \textbf{数据时效分层}：5000--9500 实时，10000 两天前，11500+ 半年前
  \item \textbf{值分布特征}：全等值=平衡三相设备（断路器），混合值=多通道设备（后备保护）
\end{enumerate}

\subsection{不依赖 Oracle 的发现路径}
本次 Tag ID 扫描\textbf{完全绕过了 Oracle 数据库}。
仅需 pSpace 连接凭据（admin/DQYTA11\_pass）和 ReadList 接口，
即可完成全作业区点位发现。
详见《OPC DA 点位自动发现方法论》。

\section{结论与展望}

\subsection{主要成果}
\begin{enumerate}[leftmargin=*]
  \item 成功逆向 psAPISDK.dll 的 Connect→ReadList 完整调用链
  \item 确认 PS\_DATA 结构体内存布局（24字节对齐）
  \item 实现自主 Tag ID 实时数据读取（python ctypes）
  \item 完成 pSpace Tag ID 空间初步探索
  \item 发现多处安全隐患（明文凭据、无加密协议）
\end{enumerate}

\subsection{待完成工作}
\begin{enumerate}[leftmargin=*]
  \item Oracle SYS\_POINTRELATION 表 → Tag ID 映射
  \item pSpace Wire Protocol 完整帧格式（需更深层逆向）
  \item psAPI\_Real\_NewSubscribeAndRead 回调机制测试
  \item 本机 pSpace mock 服务与真实 IoMonitor 对接
\end{enumerate}

\end{document}
